% Full introduction; compile the corresponding slide entry point.
\begin{frame}
\titlepage
\end{frame}

\begin{frame}{Wat u na deze inleiding kunt}
\begin{itemize}
\item Uitleggen hoe waarnemingen bruikbare gedeelde context worden.
\item Gegevens, interpretatie en operationele actie onderscheiden.
\item Bewust ontvangers en een uitwisselingsmethode kiezen.
\item Beschrijven hoe feedback informatie bruikbaar houdt.
\end{itemize}
\note{Reken op 35--40 minuten en een besprekingsoefening van tien minuten. Dit is de volledige inleiding; de korte variant duurt 15--20 minuten. Een account of API-sleutel is niet nodig. Volg één fictieve casus door het gegevensmodel, de toegangsinstellingen, de uitwisseling en de beoordeling. Praktische configuratie hoort in de vervolgcursussen.}
\end{frame}

\begin{frame}{Eén casus, twee teams}
\begin{itemize}
\item Team A onderzoekt een verdacht aanmeldbericht.
\item Het registreert \SharingCase{} en \texttt{192.0.2.25}.
\item Team B zag hetzelfde domein in een ander rapport.
\item Wat kunnen ze hergebruiken? Wat moeten ze nog vaststellen?
\end{itemize}
\note{Alle gegevens in deze casus zijn fictief. Het .example-domein en het bereik 192.0.2.0/24 zijn gereserveerd voor documentatie en zijn geen kwaadaardige indicatoren. Vraag wat Team B nodig heeft: bron, waarnemingstijd, onderbouwing en onzekerheid. Alleen een gedeeld domein bewijst geen gemeenschappelijke actor of campagne. Houd dit onderscheid tijdens de les vast.}
\end{frame}

\begin{frame}{MISP ondersteunt mensen en hulpmiddelen}
\begin{itemize}
\item Een opensourceplatform voor gestructureerde informatie-uitwisseling.
\item Waarnemingen verzamelen van analisten, tools en partners.
\item Gegevens verbinden en beoordelen met behoud van context.
\item Bruikbare invoer voor onderzoek, detectie en respons leveren.
\end{itemize}
\note{MISP ondersteunt informatie over incidenten, dreigingen, fraude en andere veiligheidsvragen. Gemeenschappen gebruiken verschillende delen van het model. Normalisatie, correlatie en verrijking helpen bij de analyse; ze valideren niet elke bewering. Het platform creëert geen vertrouwen en bepaalt niet voor iedere ontvanger de juiste reactie. Leg eerst het doel uit en daarna de functienamen.}
\end{frame}

\begin{frame}{Van waarnemingen naar gedeelde context}
\SharedIntroDiagram{sharing-flow}
\note{Volg de doorgetrokken pijlen via structuur, beoordeling en deelinstellingen naar gebruik. De stippellijn brengt feedback terug. Dit is een conceptueel proces en geen automatische uitvoeringsketen. Feeds en modules vragen configuratie en bronbeoordeling; niet elke invoer is vertrouwd en niet elke uitvoer wordt gepubliceerd. Rechten en ingestelde uitwisselingspaden begrenzen het delen.}
\end{frame}

\begin{frame}{Een klein model voor een bruikbaar record}
\SharingDiagram{data-model}
\note{Een event groepeert informatie die inhoudelijk samenhangt. Losse attributen en attributen in objecten zijn alternatieve representaties; dupliceer geen waarden alleen om elk vak te vullen. Een domain-ip-object kan een onderbouwde relatie tussen domein en IP groeperen. Objectreferenties kunnen objecten of attributen verbinden. Dit vereenvoudigde diagram is geen volledig databaseschema en legt geen exacte nesting van tags of sightings vast.}
\end{frame}

\begin{frame}{Beschrijf de waarneming vóór u handelt}
\begin{itemize}
\item Type: de vorm van een waarde, zoals \texttt{domain} of \texttt{ip-dst}.
\item Categorie: een brede indeling naar context.
\item Commentaar, bron en tijd: wat onderbouwt de waarneming?
\item \texttt{to\_ids}: selectie voor uitvoer, geen bewijs.
\end{itemize}
\note{Een attribuut kan context bieden zonder geschikt te zijn voor operationele detectie. Het type beschrijft de representatie; de categorie geeft een breed doel en vervangt de onderbouwing niet. De to-IDS-instelling wordt gebruikt door relevante exports en filters; sommige API-query's halen ook andere attributen op. Een ingeschakelde vlag maakt blokkeren niet veilig. Behandel geïmporteerde attributen niet automatisch als bevestigd.}
\end{frame}

\begin{frame}{Voeg classificatie en herbruikbare kennis toe}
\begin{itemize}
\item Taxonomietags geven gedeelde classificatietermen.
\item Galaxy-clusters voegen kennis en relaties toe.
\item ATT\&CK-technieken beschrijven onderbouwd gedrag.
\item Houd waarneming, beoordeling en onzekerheid uit elkaar.
\end{itemize}
\note{Tags kunnen aan events of attributen worden gekoppeld; lokale koppelingen hebben een andere deelreikwijdte. Galaxy-clusters bieden herbruikbare context over bijvoorbeeld actoren, malware en technieken, met relaties tussen clusters. Een actor- of ATT\&CK-label is een analytische bewering, geen door MISP gemaakt bewijs. Vermeld passende bronnen en zekerheid. De Galaxy-cursus toont de actuele interface.}
\end{frame}

\begin{frame}{Een correlatie stelt een vraag}
\SharingDiagram{correlation}
\note{Twee rapporten bevatten dezelfde domeinnaam. Werkelijke correlatie hangt af van configuratie, typen, ingeschakelde correlatie en zichtbare gegevens. Vergelijk waarnemingstijden, bronnen en gerelateerde feiten. Gedeelde hosting, hergebruikte infrastructuur, normale diensten of fouten kunnen overeenkomsten verklaren. Een overeenkomst is een onderzoeksspoor en bevestigt geen actor of campagne.}
\end{frame}

\begin{frame}{Eén record kan verschillende acties ondersteunen}
\begin{itemize}
\item Onderzoeken: rapporten vergelijken en een hypothese toetsen.
\item Detecteren of hunten: relevante waarnemingen zoeken.
\item Blokkeren: beleid toepassen en de gevolgen beoordelen.
\item Beslissing vastleggen en aangeven hoe die wordt herzien.
\end{itemize}
\note{Fout-positieven en verouderde informatie hebben andere kosten bij retrospectief zoeken dan bij automatisch blokkeren. Beoordeel bron, context, actualiteit, onzekerheid en bestemming vóór gebruik. Een gedeeld record kan nuttig zijn zonder geschikt te zijn voor blokkering. Vraag welke actie Team B nu kan onderbouwen en welk aanvullend bewijs die keuze verandert.}
\end{frame}

\begin{frame}{Kies ontvangers met de MISP-instellingen}
\small
\begin{tabular}{@{}p{.37\textwidth}p{.58\textwidth}@{}}
\toprule
Distributiekeuze & Beoogde reikwijdte \\
\midrule
Alleen uw organisatie & De organisatie die eigenaar is. \\
Alleen deze gemeenschap & De lokale MISP-gemeenschap. \\
Verbonden / alle gemeenschappen & Breder delen volgens uitwisselingsregels. \\
Sharing group & Een expliciet bepaalde groep ontvangers. \\
\bottomrule
\end{tabular}
\note{Dit zijn distributiekeuzes, geen adresboek en geen garantie dat gegevens worden verzonden. Beoordeel event, object en attribuut samen; attributen kunnen de eventdistributie erven. Neem niet aan dat beperktere gegevens automatisch een breder publiek krijgen. Zichtbaarheid hangt ook af van rollen, groepslidmaatschap en synchronisatie en filters. Beheerders kunnen bijzondere toegang hebben; organisatie-only betekent niet onzichtbaar voor de beheerder.}
\end{frame}

\begin{frame}{Omgangslabels en toegang hebben elk een taak}
\begin{itemize}
\item TLP 2.0 communiceert verwachtingen over verdere verspreiding.
\item Distributie en sharing groups regelen zichtbaarheid in MISP.
\item Beoordeel gevoelige details, doel en ontvangerverplichtingen.
\item Een label alleen dwingt geen toegestane doelgroep af.
\end{itemize}
\note{Gebruik actuele TLP 2.0-namen, waaronder TLP:CLEAR voor onbeperkt delen volgens TLP. TLP:AMBER en TLP:AMBER+STRICT hebben andere verwachtingen over verdere verspreiding; raadpleeg FIRST voor de regels. Leid geen juridische toestemming af uit een tag. Een publieke bron maakt niet elk verrijkt of gecombineerd record publiek. Stem het beleid af met de verantwoordelijke personen.}
\end{frame}

\begin{frame}{Publiceren is een bewuste stap}
\begin{enumerate}
\item Beoordeel het record, de bron en het beoogde gebruik.
\item Controleer distributie en tags op de relevante niveaus.
\item Bevestig de ingestelde sync-, feed- of exportbestemming.
\item Publiceer of wissel uit volgens het afgesproken beleid.
\end{enumerate}
\note{Opslaan, publiceren en extern uitwisselen zijn verschillende handelingen. Publicatie kan ingestelde acties activeren; paden hebben eigen rechten en filters. Ongepubliceerd betekent niet geheim: bevoegde gebruikers kunnen een record nog zien of via de API ophalen. Begin in de les met fictieve, tijdelijke gegevens, organisatie-only en ongepubliceerd. Controleer ingestelde afnemers vóór een demonstratie op een live-instantie.}
\end{frame}

\begin{frame}{Instanties wisselen geselecteerde informatie uit}
\SharingDiagram{exchange}
\note{Pijlen geven ingestelde, door beleid beperkte paden weer en geen universeel vertrouwen of verplichte centrale hub. De terugweg wordt afzonderlijk ingesteld en toegestaan. Instantie C is in dit voorbeeld losgekoppeld; goedgekeurde handmatige overdracht kan na controle wel plaatsvinden. Elke instantie kan meerdere organisaties bevatten. Beoordeel distributie, sharing groups, filters en lokaal beleid aan beide kanten. De topologie bewijst niet dat de doelgroep klopt.}
\end{frame}

\begin{frame}{Kies de uitwisselingsmethode bij het doel}
\begin{itemize}
\item Synchronisatie: gerichte uitwisseling tussen MISP-instanties.
\item Feeds: gegevens ophalen bij een ingestelde leverancier.
\item API en exports: een specifieke afnemer verbinden.
\item Handmatige overdracht: bestand en bestemming beoordelen.
\end{itemize}
\note{Een opgehaalde feed is niet automatisch een tweerichtingsgemeenschap en de leverancier ontvangt mogelijk geen correcties. Beoordeel herkomst, voorwaarden, actualiteit, dekking en overlap vóór grote imports. Native MISP JSON-feeds en andere ondersteunde feeds hebben verschillende mogelijkheden. Gebruik de feed- en synchronisatiecursussen voor details. Een losgekoppeld systeem vraagt een eigen goedgekeurd overdrachtsproces; transport maakt gegevens niet vanzelf betrouwbaar.}
\end{frame}

\begin{frame}{Werk samen met duidelijke verantwoordelijkheid}
\begin{itemize}
\item Eigenaarschap geeft aan wie het event onderhoudt.
\item Proposals laten anderen attribuutwijzigingen voorstellen.
\item Extended events voegen context van een andere organisatie toe.
\item Spreek af wie correcties ontvangt en opvolgt.
\end{itemize}
\note{Dit zijn samenwerkingsopties en geen uitwisselbare rechten. Voorstellen vragen beoordeling; extended events behouden eigen eigenaarschap en zichtbaarheid. Delegatie draagt eigenaarschap over via een afgesproken proces en garandeert geen anonimiteit. Sightings zijn een ander feedbackkanaal. Leg vast wie correcties behandelt, wijzigingen communiceert en meningsverschillen afhandelt.}
\end{frame}

\begin{frame}{Kwaliteit vraagt een terugkerende beoordeling}
\SharingDiagram{quality-cycle}
\note{Leg herkomst, waarnemingstijd en onzekerheid vast; beoordeel context en omgang; gebruik passend beleid; herbeoordeel daarna op basis van feedback. Een correctie kan de interpretatie veranderen zonder de waargenomen waarde te wijzigen. Geen score bewijst dat een record waar is. Kwaliteit en doelgroep blijven ook na publicatie aandacht vragen. De pijlen beschrijven een proces met passende automatisering en verantwoordelijkheid.}
\end{frame}

\begin{frame}{Gebruik hulpmiddelen met eigen beoordeling}
\begin{itemize}
\item Warninglists signaleren mogelijk misleidende waarden.
\item Verrijking voegt context uit gekozen bronnen toe.
\item Vergelijk herkomst en houd onzekerheid zichtbaar.
\item Controleer wat een externe module naar de leverancier stuurt.
\end{itemize}
\note{Een warninglist-hit is geen universeel oordeel dat iets fout-positief is: publieke DNS of gedeelde infrastructuur kan relevant zijn. Verrijking kan nuttige feiten toevoegen of een onbevestigde claim herhalen. Beoordeel resultaten vóór opslaan, vooral uitspraken over actoren en campagnes. Sommige modules benaderen externe diensten; controleer verzonden waarden, sleutels en beleid vóór gebruik met gevoelige gegevens. Gedrag hangt af van de configuratie.}
\end{frame}

\begin{frame}{Sightings geven verschillende soorten feedback}
\begin{itemize}
\item Sighting: een waarneming van het datapunt.
\item False positive: feedback over een ongeschikte interpretatie.
\item Expiration: feedback dat de relevantie is geëindigd.
\item Voeg bron en tijd toe; spreek af hoe feedback wordt gebruikt.
\end{itemize}
\note{MISP onderscheidt type 0 voor sighting, type 1 voor false-positive en type 2 voor expiration. Dit zijn geen drie betrouwbaarheidsscores. Een nieuwe sighting bewijst niet dat een domein kwaadaardig is of dat elk gekoppeld rapport klopt. Geen sightings bewijst geen afwezigheid van activiteit. Zichtbaarheid en terugwegen vragen configuratie en afspraken. Beloof geen automatische verwijdering, publicatie of blokkadewijziging door een sighting.}
\end{frame}

\begin{frame}{Waarnemingstijd verschilt van recordtijd}
\SharingDiagram{timeline}
\note{De fictieve datums illustreren alleen het onderscheid. first\_seen en last\_seen beschrijven, als ze bekend zijn, een waarnemingsperiode; aanmaak- en wijzigingstijden beschrijven onderhoud. Een rapport kan later worden ingevoerd dan de beschreven activiteit. Vermeld tijdzone en bron en laat onbekende tijden onbekend. Een gewijzigde attribuuttimestamp bewijst geen nieuwe waarneming. Sightingdatums zijn een andere expliciete waarnemingsregistratie.}
\end{frame}

\begin{frame}{Ouderdom en decay helpen bij herbeoordeling}
\begin{itemize}
\item Waarnemingen hebben verschillende nuttige levensduren.
\item Een decaymodel geeft een beoordelings- of filtersignaal.
\item Beoordeel het model, sightings en het beleid van de afnemer.
\item Bewaar historische context na operationeel verval.
\end{itemize}
\note{Decay is geen universele kans op malware en geen opdracht om oude gegevens te verwijderen. Effecten op exports of acties hangen af van modellen, instellingen, filters en integraties. Verlopen operationele relevantie kan samengaan met historische onderzoekswaarde. Gebruik de decaycursussen voor configuratie en simulatie in plaats van hier oude knoplocaties te onderwijzen.}
\end{frame}

\begin{frame}{Automatiseer een beoordeelde beslissing}
\begin{itemize}
\item Gebruik de REST API en PyMISP voor afgebakende processen.
\item Filter voor de afnemer, het doel en de doelgroep.
\item Test met fictieve gegevens en inspecteer de uitvoer.
\item Breng feedback terug en behandel fouten expliciet.
\end{itemize}
\note{Een integratie moet vastleggen wie wat mag lezen of wijzigen, welke gegevens nodig zijn, welke filters gelden en wat er bij herhaling of fouten gebeurt. Gebruik een beperkt serviceaccount en bescherm de API-sleutel. Deze inleiding laat versiegevoelige payloads en methodesignaturen bewust weg; de vernieuwde REST API- en PyMISP-cursussen behandelen die. Bewaar herkomst en toegangsbeleid in de automatisering.}
\end{frame}

\begin{frame}{Een gemeenschap heeft werkafspraken nodig}
\begin{itemize}
\item Spreek doel, deelnemers en bijdragen af.
\item Bepaal omgang, correctie en escalatie.
\item Begin met een kleine vertrouwde groep en nuttige records.
\item Herzie afspraken als mensen en toepassingen veranderen.
\end{itemize}
\note{Een gemeenschap kan een intern team, sectorgroep, CSIRT-netwerk, fraudepartnerschap of thematische samenwerking zijn. Een instantie is een implementatie; een sharing group is een doelgroep; beide verschillen van de sociale gemeenschap. Vertrouwen groeit uit relaties en werkwijzen. Software ondersteunt de afspraken maar levert niet hun juridische basis, motivatie of verantwoording.}
\end{frame}

\begin{frame}{Maak van een obstakel een eerste bijdrage}
\small
\begin{tabular}{@{}p{.28\textwidth}p{.68\textwidth}@{}}
\toprule
Obstakel & Een nuttige eerste stap \\
\midrule
Weinig tijd & Deel een kleine waarneming met een bron. \\
Onzekerheid & Benoem wat bekend en corrigeerbaar is. \\
Andere tools & Spreek termen en ondersteunde formaten af. \\
Verplichtingen & Beoordeel doel en ontvangers met verantwoordelijken. \\
\bottomrule
\end{tabular}
\note{Doe technische, organisatorische of juridische beperkingen niet af als denkbeeldig. Werk met de verantwoordelijke mensen. Een bronnotitie, correctie of contextuele link kan nuttiger zijn dan nog een grote indicatorenlijst. Goede voorbeelden en eenvoudige feedback helpen deelname. Vraag om één haalbare bijdrage onder het eigen beleid, zonder gevoelige gegevens in de klas op te vragen.}
\end{frame}

\begin{frame}{Het project begon met een praktische behoefte}
\begin{itemize}
\item In 2012 wilden malwareanalisten dubbel werk voorkomen.
\item Het vroege platform van Christophe Vandeplas werd MISP.
\item CIRCL en een bredere gemeenschap verbeteren het project.
\item Bijdragen omvatten code, gedeelde kennis en opleiding.
\end{itemize}
\note{Behoud de oorsprong uit het eerdere lesmateriaal: een malwarewerkgroep ontdekte dubbel werk en gebruikersfeedback vormde het platform. CIRCL speelt al lang een rol in ontwikkeling en gemeenschap. Verwijs naar actuele CIRCL-informatie in plaats van ministeriële mandaten of organisatiestructuren vast te zetten. Vermijd ongedateerde gebruikersaantallen en beperk de huidige gemeenschappen niet tot de eerste malwaretoepassing.}
\end{frame}

\begin{frame}{Het MISP-project heeft vier aanvullende delen}
\SharedIntroDiagram{overview}
\note{Software, kennis, standaarden en gemeenschappen behouden het bruikbare vierdelige idee uit het oorspronkelijke diagram. Categorieën overlappen en vormen geen hiërarchie. MISP, PyMISP en optionele modules ondersteunen gebruik; taxonomieën, Galaxies en warninglists bieden kennis; formaten, begrippen en sjablonen helpen uitwisseling; mensen leveren werkwijzen en feedback. Niet elk onderdeel hoeft geïnstalleerd te worden.}
\end{frame}

\begin{frame}{De gemeenschap verbetert het project}
\SharedIntroDiagram{community}
\note{Dit vervangt het oude governanceplaatje door een bewerkbare conceptuele bijdragecyclus. Gebruikers brengen behoeften; bijdragers stellen code, gegevens, documentatie en opleiding voor; beheerders beoordelen en nemen releasebeslissingen; verbeteringen keren terug naar gebruikers. Mensen kunnen meerdere rollen hebben. Het is geen formeel juridisch of organisatorisch gezagsschema. Volg de echte bijdrageprocedure van de betreffende repository.}
\end{frame}

\begin{frame}{Oefening: maak de casus bruikbaar voor Team B}
\begin{enumerate}
\item Ontwerp een event met het fictieve domein en IP-adres.
\item Leg bron, tijd, relatie en onzekerheid uit.
\item Kies beperkte ontvangers en een onderbouwde vervolgactie.
\item Beschrijf hoe een correctie Team B bereikt.
\end{enumerate}
\note{Geef tien minuten in tweetallen of als groepsgesprek. Papier is voldoende. Gebruik in een toegestaan lab alleen tijdelijke fictieve gegevens, houd het event organisatie-only en ongepubliceerd en controleer ingestelde afnemers. Publiceer of blokkeer de documentatiewaarden niet. Vraag onderscheid tussen waarneming en beoordeling en een uitleg van de gekozen ontvangers.}
\end{frame}

\begin{frame}{Nabespreking: leg de beslissingen uit}
\begin{itemize}
\item Het gedeelde domein is een spoor, geen attributie.
\item Een object groepeert een door de bron onderbouwde relatie.
\item Omgangslabels en zichtbaarheid vragen beide controle.
\item Een zoekactie kan eerder verantwoord zijn dan blokkeren.
\item Een correctie heeft een eigenaar en een terugweg nodig.
\end{itemize}
\note{Zonder gemeenschapsbeleid en bewijs is er geen unieke juiste sharing group of blokkeerbeslissing. Een goed antwoord bewaart herkomst, benoemt ontbrekende feiten en beperkt ontvangers bewust. Vraag wat verandert als het IP gedeelde infrastructuur is, de bron haar claim intrekt of de activiteit maanden oud is. Bespreek besliskwaliteit in plaats van een actor te raden.}
\end{frame}

\begin{frame}{Ga verder met de taak die u nodig hebt}
\begin{itemize}
\item Dagelijkse analyse: gebruik, tags en objecten.
\item Context en kwaliteit: Galaxies en indicatorverval.
\item Integratie: REST API, PyMISP en feeds.
\item Gemeenschapsbeheer: administratie en deelpraktijken.
\end{itemize}
\note{De README verwijst direct naar actuele cursussen en bronnen. Kies vervolgcursussen op basis van de rol van de deelnemers in plaats van elke productfunctie te doorlopen. De volledige inleiding geeft een gedeelde woordenschat en redeneerproces. Trainers kunnen TikZ-diagrammen eenmaal aanpassen voor beide talen of als losse vectorbestanden exporteren.}
\end{frame}

\begin{frame}{Deel context, kies bewust en blijf verbeteren}
\begin{itemize}
\item Structureer waarnemingen met behoud van betekenis.
\item Kies ontvangers en acties volgens afgesproken beleid.
\item Gebruik feedback voor correctie en herbeoordeling.
\item Draag ook werkwijzen en kennis bij.
\end{itemize}
\note{Vraag de deelnemers de casus in één minuut samen te vatten: wat is waargenomen, wat ondersteunt dat, wie mag het ontvangen, welke actie volgt en hoe kan het worden gecorrigeerd? Het doel is nuttige informatie-uitwisseling met MISP en geen functiecatalogus. Een kleine, goed uitgelegde bijdrage kan een ander team dubbel werk of een slechte beslissing besparen.}
\end{frame}

